\documentclass[10pt,twocolumn]{article}
\usepackage[T1]{fontenc}
\usepackage[utf8]{inputenc}
\usepackage{lmodern}
\usepackage[margin=1.8cm,top=2.4cm,bottom=2cm,columnsep=0.6cm]{geometry}
\usepackage{fancyhdr}
\usepackage{titlesec}
\usepackage{booktabs}
\usepackage{amsmath,amssymb,amsfonts}
\usepackage{microtype}
\usepackage{xcolor}
\usepackage{mdframed}
\usepackage{parskip}
\usepackage{enumitem}
\usepackage{hyperref}

\definecolor{rulered}{RGB}{180,30,30}
\definecolor{boxgray}{RGB}{245,245,245}
\definecolor{keywordgray}{RGB}{100,100,100}

\pagestyle{fancy}
\fancyhf{}
\fancyhead[L]{\textsc{\small Claude Mag}}
\fancyhead[C]{\small\textit{Machine Learning Research Digest}}
\fancyhead[R]{\small 2026-10-21}
\fancyfoot[C]{\small\thepage}
\renewcommand{\headrulewidth}{0.8pt}

\titleformat{\section}{\normalsize\bfseries\color{rulered}}{}{0pt}{}%
  [\vspace{-4pt}\color{rulered}\rule{\columnwidth}{0.4pt}]
\titlespacing{\section}{0pt}{8pt}{4pt}

\hypersetup{colorlinks=true,urlcolor=rulered,linkcolor=black}

\begin{document}
\setlength{\parindent}{0pt}
\setlength{\parskip}{4pt}

{\small\color{keywordgray}\textbf{Keywords:}
  computer-use agents \textbullet{}
  environment synthesis \textbullet{}
  synthetic training data \textbullet{}
  GUI agents}\par\vspace{4pt}

{\LARGE\bfseries\raggedright
  Agent World Model in the Wild: Computer-Use Agent
  Environment Synthesis}\par\vspace{4pt}

{\large\itshape\raggedright
  Persona-Grounded WildAWM Synthesises 341 Long-Horizon Computer-Use
  Environments; Qwen Models Trained There Generalise to 8 OOD Benchmarks}\par\vspace{6pt}

{\small\textbf{By} Zhaoyang Wang, Qianhui Wu, Huaxiu Yao, et al.
  (Microsoft Research \& UNC Chapel Hill)}\par\vspace{2pt}

{\small\color{keywordgray}arXiv:2610.26828 \textbullet{}
  October 2026}
\par\vspace{2pt}\noindent{\color{rulered}\rule{\columnwidth}{0.6pt}}\vspace{6pt}

Training general-purpose computer-use agents requires controllable
environments with verifiable ground truth---a resource that is scarce,
expensive to build manually, and impossible to extract from screen
recordings without execution-based verification. WildAWM synthesises such
environments by grounding them in synthetic user personas: it generates
executable desktop worlds with persona-appropriate websites, MCP services,
and files, then synthesises long-horizon cross-application tasks and filters
them through execution-based quality gates combining heterogeneous agent
rollouts, cross-judge consistency, and sycophancy-probing counterfactuals.
The resulting 341 environment-task pairs, which GPT-6-Astra solves at only
60\%, transfer well: Qwen models trained exclusively on them gain substantially
across 8 out-of-distribution computer-use benchmarks.

\begin{mdframed}[backgroundcolor=boxgray,linecolor=rulered,linewidth=1pt,
    innerleftmargin=6pt,innerrightmargin=6pt,
    innertopmargin=6pt,innerbottommargin=6pt]
\textbf{\small AT A GLANCE}\par\vspace{4pt}
\begin{description}[leftmargin=0pt,labelsep=4pt,itemsep=2pt,parsep=0pt,
    font=\small\bfseries]
  \item[Problem:] Realistic, verifiable computer-use training environments
    are scarce; scripted generators produce tasks too simple to transfer.
  \item[Why it matters:] Frontier models are approaching human performance
    on curated benchmarks; progress now requires harder, more diverse
    training data.
  \item[Method:] Persona-grounded environment synthesis with three-tier
    quality gate: heterogeneous rollouts, cross-judge consistency, and
    counterfactual sycophancy probes.
  \item[Result:] 341 long-horizon pairs; GPT-6-Astra at 60\%; Qwen models
    trained on synthetic data generalise to 8 OOD benchmarks.
  \item[Evidence:] OOD evaluation on web, desktop, and computer-use tasks;
    quality gate ablation confirms each tier's contribution.
\end{description}
\end{mdframed}

\section{The Big Picture}

The ceiling for computer-use agents is no longer model architecture or
training objective---it is training data. Frontier models like GPT-6-Astra
now score above 90\% on the original OSWorld and ScreenSpot benchmarks;
further gains require harder, more diverse, and more verifiable training
environments.

Building such environments by hand does not scale. WildAWM takes a generative
approach: start from a synthetic user persona (a profile describing the
person's job, preferences, and daily computer activities), derive an
executable personal-computer environment from that persona, and synthesise
tasks that are grounded in and verifiable against that world.

The result is a data pipeline that can, in principle, generate environments
indefinitely---limited only by the diversity of the persona generator and
the quality of the downstream verifier.

\section{Why Existing Approaches Fall Short}

\textbf{Hand-curated benchmarks (OSWorld, ScreenSpot, AndroidWorld):}
Expensive to build; fixed task distribution that saturates as models
improve; evaluation sets cannot scale to training-set sizes.

\textbf{Rule-based environment generators:} Tasks are formulaic (``click
button X on website Y''); they do not resemble the cross-application,
multi-step workflows that real users perform.

\textbf{Web-crawled demonstrations:} Ground-truth completion is difficult
to verify; rewards are noisy; task length and difficulty cannot be controlled.

\textbf{Simulation replay:} Requires access to the original software
environment; does not generalise to novel configurations.

\section{Core Mechanism}

\textbf{Persona-grounded environment synthesis:} An LLM generates synthetic
user personas covering profession, daily activities, and software usage
patterns. From each persona, WildAWM constructs an executable desktop
environment: a browser populated with persona-relevant web pages, MCP
(Model Context Protocol) services simulating the APIs the persona would
use, and a local file system with persona-appropriate documents.

\textbf{Cross-agent interface (XAI) task synthesis:} Tasks that require
actions across multiple applications (e.g., ``find the flight mentioned
in your last email, book it, and create a packing list in your documents'')
are synthesised as long-horizon plans with verifiable sub-goals. These
XAI tasks mirror the cross-application coordination that distinguishes
real user workflows from benchmark tasks.

\textbf{Three-tier quality gate:}
\begin{enumerate}[itemsep=2pt,parsep=0pt]
  \item \textbf{Heterogeneous rollouts:} Multiple agent architectures
    attempt each task. Tasks trivially solved by all (too easy) or
    unsolvable by any (unverifiable) are filtered.
  \item \textbf{Cross-judge consistency:} Multiple judge models evaluate
    completion; pairs with inconsistent verdicts are rejected.
  \item \textbf{Counterfactual sycophancy probes:} Judges are presented
    with deliberately incorrect outcome feedback; judges that agree with
    false outcomes are excluded.
\end{enumerate}

\section{Technical Formulation}

Let $u \sim P_{\text{persona}}$ be a synthetic user persona,
$e \sim P_{\text{env}}(\cdot \mid u)$ the derived environment, and
$\tau \sim P_{\text{task}}(\cdot \mid e)$ a synthesised task.

The quality gate passes $(e, \tau)$ if:
\[
Q(e, \tau) = \mathbf{1}\!\left[\text{ConsistencyScore}(e, \tau) > \delta_c\right]
  \cdot \mathbf{1}\!\left[\text{RobustnessScore}(e, \tau) > \delta_r\right]
\]

The \textit{robustness score} exploits LLM sycophancy to stress-test verifiers:
\[
\text{RobustnessScore}(e, \tau)
  = 1 - \frac{1}{|C|}\sum_{c \in C} \mathbf{1}\!\left[
    \text{Judge}(e, \tau, c_{\text{false}}) = \texttt{correct}\right]
\]
where $c_{\text{false}}$ is a deliberately incorrect outcome label; a
judge that accepts it is flagged as sycophantic and excluded from the
verifier pool.

A Qwen model $\pi_\theta$ is then trained with behavioural cloning on
the curated pairs followed by RL with the verifier as the reward model:
\[
\mathcal{J}(\theta) = \mathbb{E}_{(e,\tau)\sim\mathcal{D}_{\text{curated}}}
  \!\left[\mathbb{E}_{a\sim\pi_\theta(\cdot|e,\tau)}\!\left[
    V(e, \tau, a)\right]\right]
\]
where $V$ is the quality-gated verifier score.

\section{Learning or Inference Procedure}

\begin{enumerate}[itemsep=2pt,parsep=0pt]
  \item Sample personas; generate environments via LLM; synthesise XAI tasks.
  \item Apply the three-tier quality gate; retain pairs that pass all checks.
  \item Train Qwen models on curated pairs with behavioural cloning, then
    RL with the verified reward.
  \item Evaluate on held-out OOD benchmarks: no fine-tuning on any target
    benchmark.
\end{enumerate}

\section{What the Guarantee Says}

The counterfactual sycophancy probe provides a robustness guarantee over
the verifier pool: any judge that accepts a known-false outcome is excluded,
bounding the false-positive rate of the verification signal. The consistency
check across heterogeneous judges bounds the variance of the reward signal.
Together, the two gates ensure that the training signal is both accurate and
stable, which is the precondition for effective RL in the synthetic environment.

\section{Experimental Findings}

\begin{itemize}[itemsep=2pt,parsep=0pt]
  \item \textbf{Dataset:} 341 long-horizon task-environment pairs, each
    requiring up to 100+ agent steps. GPT-6-Astra solves only 60\%,
    confirming non-trivial difficulty.
  \item \textbf{OOD generalisation:} Qwen models trained exclusively on
    WildAWM environments gain substantially across 8 OOD benchmarks
    spanning web browsing, email management, file management, coding
    assistance, and general computer-use tasks.
  \item \textbf{Cross-application advantage:} Gains are largest on
    multi-application coordination tasks, the XAI task category WildAWM
    specialises in.
  \item \textbf{Quality gate ablation:} Removing any single tier of the
    gate degrades training-set quality (measured by human evaluation) and
    downstream performance.
\end{itemize}

\section{Ablations and Interpretation}

\textbf{Without persona grounding (random environments):} Task-environment
pairs are less realistic; downstream benchmark performance drops, especially
on tasks requiring knowledge of realistic file and browser content.

\textbf{Without counterfactual probes:} Some judges are sycophantic and
accept false completions; reward signal becomes noisier; downstream
performance degrades.

\textbf{Without cross-judge consistency:} High-variance verification
labels slow convergence and reduce final performance.

\textbf{Why synthetic environments transfer:} Persona-grounded synthesis
produces task distributions closely aligned with real user workflows; the
XAI structure requires the same cross-application reasoning as real tasks,
rather than isolated widget interactions that do not compose.

\section*{Reference}

Zhaoyang Wang, Qianhui Wu, Huaxiu Yao, et al.
\textbf{Agent World Model in the Wild: Computer-Use Agent Environment Synthesis}.
arXiv:2610.26828, October 2026.
\url{https://arxiv.org/abs/2610.26828}

\end{document}
